\section{Introduction}
\label{sec:intro}

The existence of programs like the \Gls{GAIN} \cite{doe_gain_2025}, the
\Gls{ARDP} \cite{doe_ardp_2025}, and the \Gls{RPP} \cite{doe_rpp_2025} show
that there is serious private sector, academic, and governmental interest in
advanced reactor designs. Computational tools capable of supporting the design,
engineering, and analysis of these reactors are needed. Much of the
engineering, safety analysis, and design work required for nuclear power
reactors, in particular those pursuing a first design license, must account for
situations outside of steady state operation. Deterministic methods are
typically used to handle this, however there is growing interest in utilizing
time-dependent Monte Carlo as a high-fidelity complement to time-dependent
deterministic methods \cite{faucher_new_2018, jang_advances_2023,
    romero-barrientos_td_2022, legrady_pb_2020, shaw_hybrid_2025,
molnar_gpu_2019}.

Time-dependent Monte Carlo simulations are computationally expensive, so
variance reduction techniques are typically used to improve simulation
performance. Time-dependent Cooper-Larsen weight windows have been shown to
greatly increase the average \Gls{FOM} of a time-dependent simulation by five
orders of magnitude compared to an analog simulation \cite{northrop_study_2024}.
\Gls{FW-CADIS}, the ``gold standard'' \cite{munk_review_2019} of automatic
variance reduction methods, has not been applied to time-dependent Monte Carlo
in any currently published literature.  \Gls{FW-CADIS} weight windows have been
shown to increase the \Gls{FOM} of a Monte Carlo up to two orders of magnitude
over Cooper-Larsen weight windows, depending on the problem
\cite{peplow_comparison_2012}. It is expected that this increase in \Gls{FOM}
will be reflected in time-dependent problems as well.

A time-dependent \Gls{FW-CADIS} implementation requires both a time-dependent
Monte Carlo method {\it and} a time-dependent deterministic method to obtain the
adjoint flux solution at each time step. \OpenMC, an open-source neutron
particle transport code, recently gained a deterministic calculation mode based
on the Random Ray Method \cite{tramm_random_2024}. The Random Ray Method is a
form of the \gls{MOC} that uses random sampling to generate quadrature. This has
several attractive properties for parallelization and memory usage that are
highly beneficial if used for time-dependent simulations. The Random Ray
implementation in \OpenMC already supports \Gls{FW-CADIS} for steady state
problems. This has the benefit that the same geometry files can be used for the
Monte Carlo and Random Ray simulations, provided the appropriate multi-group
energy treatment is specified for the Random Ray simulations.

